% Fisher Curvature Conjecture Note — layout consistent with LLM_Grammar_Fisher_Preprint.tex
% Build: pdflatex Fisher_Curvature_Conjecture_Note.tex (twice) from docs/
\documentclass[11pt,a4paper]{article}
\usepackage[T1]{fontenc}
\usepackage[utf8]{inputenc}
\usepackage{lmodern}
\usepackage{geometry}
\geometry{margin=1in}
\usepackage{amsmath,amssymb}
\usepackage{microtype}
\usepackage{booktabs}
\usepackage{fancyhdr}
\usepackage{lastpage}
\usepackage[hidelinks]{hyperref}
\usepackage{csquotes}

\setlength{\parskip}{0.45em}
\setlength{\parindent}{0pt}

\pagestyle{fancy}
\fancyhf{}
\fancyhead[L]{\small Curvature at the Fisher Threshold --- Research Note (April 2026)}
\fancyhead[R]{\small \thepage\ / \pageref{LastPage}}
\renewcommand{\headrulewidth}{0.4pt}

\begin{document}

\begin{center}
{\large Research Note --- April 2026}\\[0.8em]
{\LARGE\bfseries Curvature Structure at the Fisher Information Threshold:\\[0.25em]
Numerical Evidence from Quantum Processors and Language Models}\\[0.9em]
{\large Continuation of \textit{Grammar Fingerprinting} (\href{https://doi.org/10.5281/zenodo.19158088}{doi:10.5281/zenodo.19158088}),\\
\textit{Fisher Information Threshold Study} (\href{https://doi.org/10.5281/zenodo.19391582}{doi:10.5281/zenodo.19391582})}\\[0.5em]
Daniel Csapl\'ar\\
Independent Researcher, Kazincbarcika, Hungary\\
ORCID: \href{https://orcid.org/0009-0000-7362-7232}{0009-0000-7362-7232}\\
April 2026
\end{center}

\vspace{0.6em}
\noindent\textbf{Abstract}\\[0.3em]
We report numerical evidence for a curvature structure at the Fisher information threshold $N^\star$ in Grammar Fingerprinting transition matrices. Using second-derivative analysis of the Fisher trace $\mathrm{Tr}(F)$ as a function of data length $N$, we show that the global maximum of $|d^2\mathrm{Tr}/dN^2|$ coincides with $N^\star$ (the Fisher trace maximum) in 89\% of 28 Google Sycamore quantum processor readout configurations (within one grid step). In LLM entropy time series (Qwen2.5-Coder:7B, 10{,}000 tokens), the coincidence holds for simpler regimes (factual), while more complex regimes (mathematical) exhibit a separation between an early curvature peak and a later trace maximum---consistent with multi-scale temporal structure. We propose that the Fisher trace maximum and the curvature maximum define two complementary scales: a structure-onset threshold (where the sharpest information gain occurs) and an information ceiling (where the learned grammar saturates). In simple systems these coincide; in complex systems they separate. We conjecture that this curvature structure reflects a geometric phase transition on the statistical manifold of transition matrices, and note the structural parallel to the macroscopicity threshold in the Page--Wootters mechanism for emergent time. The formal derivation of this connection remains an open problem.

\vspace{0.6em}
\textbf{1.\ The Conjecture}

We state the conjecture in three parts.

\textbf{Part 1 (Curvature--threshold coincidence):} For structured sequential processes with a single dominant temporal scale, the global maximum of $|d^2\mathrm{Tr}(F)/dN^2|$ coincides with the Fisher trace maximum $N^\star$ on the data-length grid.

\textbf{Part 2 (Scale separation in complex systems):} For processes with multiple temporal scales, the global curvature maximum $N_{\mathrm{curv}}$ and the Fisher trace maximum $N^\star$ separate. $N_{\mathrm{curv}}$ marks the onset of structure acquisition (sharpest information gain); $N^\star$ marks the saturation of the dominant grammar. The ratio $N^\star/N_{\mathrm{curv}}$ reflects the system's structural complexity.

\textbf{Part 3 (Geometric phase transition):} The curvature structure at $N^\star$ (or $N_{\mathrm{curv}}$) corresponds to a change in the Riemannian curvature of the statistical manifold of row-stochastic transition matrices, parameterized by data length. This change is analogous to the macroscopicity threshold in the Page--Wootters (PaW) mechanism, where classical time emerges from quantum entanglement above a critical system size.

Parts 1 and 2 are supported by numerical evidence presented below. Part 3 is a formal conjecture requiring analytical derivation.

\vspace{0.5em}
\textbf{2.\ Numerical Evidence: Google Sycamore}

\textbf{2.1 Data and method.} Fisher trace curves $\mathrm{Tr}(F)$ vs.\ data length $N$ were computed for all 28 readout configurations of the Google Sycamore quantum supremacy dataset (Arute et al., 2019), using the Grammar Fingerprinting pipeline (SAX alphabet 7, LSTM hidden 32, seq\_len 16, 50 epochs). The grid spans $N=500$ to 50{,}000 in 21 steps. For each curve, the second derivative $d^2\mathrm{Tr}/dN^2$ was computed via natural cubic spline interpolation, and the global maximum of $|d^2\mathrm{Tr}/dN^2|$ was identified.

\textbf{2.2 Results.}
\begin{center}
\begin{tabular}{@{}lr@{}}
\toprule
Criterion & Result \\
\midrule
Exact coincidence ($\Delta=0$) & 19 / 28 (67.9\%) \\
Within one grid step ($\Delta\le 1000$) & 25 / 28 (89.3\%) \\
Within two grid steps ($\Delta\le 2000$) & 26 / 28 (92.9\%) \\
Among high-curvature readouts ($|d^2|\ge 100$) & 9 / 9 within $\Delta\le 1000$ \\
Outliers (46q, 49q) & $|d^2|<0.1$ --- flat curvature, peak location unreliable \\
\bottomrule
\end{tabular}
\end{center}

The median $\Delta$ is 0; the mean is inflated by two outliers with essentially flat second derivatives. Where the curvature is well-defined ($|d^2|\ge 100$), the coincidence is complete.

\textbf{2.3 Interpretation.} On Google Sycamore, each readout configuration represents a single physical topology with one dominant temporal scale. The Fisher trace rises, peaks at $N^\star$, and declines. The curvature maximum coincides with this peak, consistent with Part 1 of the conjecture: in a single-scale system, the sharpest information gain and the information ceiling occur at the same data length.

\vspace{0.5em}
\textbf{3.\ Numerical Evidence: LLM Entropy Series}

\textbf{3.1 Data and method.} Fisher trace curves were computed for entropy time series from Qwen2.5-Coder:7B (10{,}000 tokens, LLM seed 42) under factual and mathematical discourse regimes, each with three conditions: real signal (\emph{none}) and two controls (\emph{shuffled}, \emph{random\_uniform}). The Fisher pipeline matches the primary LLM study in the repository: SAX alphabet 7, LSTM hidden 32, seq\_len 16, 50 epochs, with \textbf{grammar (LSTM) seed 2} for the reported Fisher-vs-$N$ curves (same convention as \texttt{llm\_entropy\_fisher\_summary} and related outputs). The same spline-based second derivative analysis was applied (\texttt{curvature\_fisher\_trace\_llm.py}).

\textbf{3.2 Results.} Table~1 lists $N^\star=\arg\max_N \mathrm{Tr}(F)$, $N_{\mathrm{curv}}=\arg\max_N |d^2\mathrm{Tr}/dN^2|$, and $\Delta=|N_{\mathrm{curv}}-N^\star|$ for all six curves.

\begin{table}[htbp]
\centering
\caption{LLM seed 42: factual and mathematical regimes (grammar seed 2).}
\begin{tabular}{@{}llrrr@{}}
\toprule
Regime & Condition & $N^\star$ & $N_{\mathrm{curv}}$ & $\Delta$ \\
\midrule
Factual & none & 384 & 384 & 0 \\
Factual & shuffled & 224 & 224 & 0 \\
Factual & random\_uniform & 32 & 48 & 16 \\
Mathematical & none & 1536 & 112 & 1424 \\
Mathematical & shuffled & 48 & 48 & 0 \\
Mathematical & random\_uniform & 64 & 64 & 0 \\
\bottomrule
\end{tabular}
\end{table}

The factual regime (lower structural complexity in this comparison) shows exact or near coincidence. The mathematical \emph{none} condition shows clear scale separation: the curvature peak occurs early ($N_{\mathrm{curv}}=112$, $|d^2|\approx 476$) while the Fisher trace peaks much later ($N^\star=1536$). A weak local $|d^2|$ feature exists at $N=1536$ ($|d^2|\approx 11$) but is not dominant relative to the early peak.

\textbf{3.3 Interpretation.} Mathematical discourse in an LLM involves multiple structural layers: syntactic constraints, logical sequencing, symbolic conventions, and topic-level coherence. The early curvature peak ($N_{\mathrm{curv}}=112$) likely reflects the LSTM acquiring basic transition structure; the Fisher trace continues to rise as longer-range dependencies accumulate, peaking at $N^\star=1536$. This is consistent with Part 2 of the conjecture: in multi-scale systems, $N_{\mathrm{curv}}$ and $N^\star$ separate, and their ratio ($N^\star/N_{\mathrm{curv}}\approx 14$ here) may serve as a complexity indicator.

\vspace{0.5em}
\textbf{4.\ Two Scales, One Mechanism}

We propose that the Grammar Fingerprinting pipeline, applied to any structured sequential process, defines two characteristic scales:

\textbf{$N_{\mathrm{curv}}$ (structure-onset threshold):} The data length at which the sharpest change in Fisher information occurs. This is where the learning system first ``catches'' the dominant pattern. Operationally: $\arg\max_N |d^2\mathrm{Tr}(F)/dN^2|$.

\textbf{$N^\star$ (information ceiling):} The data length at which the Fisher trace is maximal. This is where the learned grammar carries the most information about the transition structure. Operationally: $\arg\max_N \mathrm{Tr}(F)$.

In simple systems (single temporal scale), $N_{\mathrm{curv}}\approx N^\star$. In complex systems (multiple temporal scales), $N_{\mathrm{curv}}<N^\star$, and the ratio $N^\star/N_{\mathrm{curv}}$ characterizes the system's structural depth.

This framework is consistent with the transfer model from the Fisher threshold study (Csapl\'ar, 2026b):
\[
N^\star = H_0 \times B_{\mathrm{backend}} \times C_{\mathrm{regime}} \times Q^{\alpha},
\]
where the threshold depends on the backend ($B$), the circuit or discourse regime ($C$), and the system size ($Q$) with exponent $\alpha\approx -0.77$. The present analysis adds that the internal structure of the transition---whether it is a single sharp event or a multi-stage process---is visible in the curvature profile.

\vspace{0.5em}
\textbf{5.\ Connection to Emergent Time (Open Problem)}

The Page--Wootters (PaW) mechanism (Page \& Wootters, 1983; Foti et al., 2021) shows that time emerges from entanglement between two non-interacting quantum subsystems when the ``clock'' subsystem reaches a macroscopicity threshold. Below this threshold, the composite system is static (Wheeler--DeWitt regime); above it, unitary evolution---and thus time---emerges.

The structural parallel to Grammar Fingerprinting is direct:

PaW: below macroscopicity threshold $\rightarrow$ static, no time. Above $\rightarrow$ time emerges.

Grammar Fingerprinting: below $N^\star$ $\rightarrow$ no detectable temporal structure (chance-level classification, Fisher trace at noise floor). Above $N^\star$ $\rightarrow$ temporal grammar emerges.

Both transitions are sharp, system-dependent, and mark the boundary between structurelessness and emergent order.

We conjecture (Part 3) that this parallel is not merely analogical. If the statistical manifold of transition matrices $T(N)$---parameterized by data length $N$---has a Riemannian structure defined by the Fisher information metric, then the curvature singularity at $N^\star$ (or $N_{\mathrm{curv}}$) may correspond formally to the quantum-to-classical crossover in the PaW framework, where the large-$N$ limit of generalized coherent states produces classical equations of motion.

The formal derivation of this correspondence---connecting the curvature of the transition-matrix manifold to the macroscopicity condition in canonical quantum gravity---remains an open problem. It requires tools from information geometry (Amari, 1985), the large-$N$ coherent state formalism (Foti et al., 2021), and the theory of statistical phase transitions on simplicial probability spaces.

This note provides the numerical motivation; the analytical bridge is left for future work.

\vspace{0.5em}
\textbf{6.\ Summary}

\begin{enumerate}\setlength\itemsep{0.25em}
\item On Google Sycamore (28 readouts), the curvature maximum $|d^2\mathrm{Tr}/dN^2|$ coincides with the Fisher trace maximum $N^\star$ in 89\% of configurations (within one grid step). Where the curvature is well-defined, the coincidence is complete.
\item On LLM entropy series, the coincidence holds for simpler discourse regimes (factual). For complex regimes (mathematical none), the curvature peak and trace maximum separate, defining two scales: a structure-onset threshold ($N_{\mathrm{curv}}$) and an information ceiling ($N^\star$).
\item We conjecture that this curvature structure reflects a geometric phase transition on the statistical manifold of learned transition matrices, with a formal parallel to the macroscopicity threshold in the Page--Wootters mechanism for emergent time.
\item The formal derivation of this connection is an open problem.
\end{enumerate}

\vspace{0.6em}
\noindent\textbf{References}

\begin{itemize}\setlength\itemsep{0.2em}
\item Amari, S. (1985). \textit{Differential-Geometrical Methods in Statistics}. Springer.
\item Arute, F., et al.\ (2019). Quantum supremacy using a programmable superconducting processor. \textit{Nature}, 574, 505--510.
\item Csapl\'ar, D. (2026a). Grammar Fingerprinting of Quantum Processor Topology. Zenodo. \href{https://doi.org/10.5281/zenodo.19158088}{doi:10.5281/zenodo.19158088}.
\item Csapl\'ar, D. (2026b). Fisher Information Threshold Study: Grammar Fingerprinting on Google Sycamore. Zenodo. \href{https://doi.org/10.5281/zenodo.19391582}{doi:10.5281/zenodo.19391582}.
\item Csapl\'ar, D. (2026c). Grammar Fingerprinting and Fisher Information Thresholds in Large Language Model Entropy Series. Zenodo. \href{https://doi.org/10.5281/zenodo.19454201}{doi:10.5281/zenodo.19454201}.
\item Foti, C., et al.\ (2021). Time and classical equations of motion from quantum entanglement via the Page and Wootters mechanism with generalized coherent states. \textit{Nature Communications}, 12, 1787.
\item Page, D.\ N.\ \& Wootters, W.\ K.\ (1983). Evolution without evolution: Dynamics described by stationary observables. \textit{Physical Review D}, 27, 2885.
\end{itemize}

\vspace{0.6em}
\noindent\textbf{Data and Code Availability}

Sycamore readout data: Dryad DOI \href{https://doi.org/10.5061/dryad.k6t1rj8}{10.5061/dryad.k6t1rj8}. Grammar Fingerprinting code and validation results: Zenodo \href{https://doi.org/10.5281/zenodo.19158088}{doi:10.5281/zenodo.19158088}. Fisher threshold study: Zenodo \href{https://doi.org/10.5281/zenodo.19391582}{doi:10.5281/zenodo.19391582}. LLM entropy study: Zenodo \href{https://doi.org/10.5281/zenodo.19454201}{doi:10.5281/zenodo.19454201}.

\begin{sloppypar}
Curvature analysis scripts (\texttt{curvature\_fisher\_trace\_sycamore.py}, \texttt{curvature\_fisher\_trace\_llm.py}), summary CSVs, and per-$N$ output files are maintained in the project repository:\\
\url{https://github.com/csaplard/QuantumCircuit_Grammar_Research}
\end{sloppypar}

\vspace{0.8em}
\begin{center}
--- End of research note ---
\end{center}

\end{document}
